% mode: ini-etex
% WEB 1447: the token saved by \afterassignment is inserted after every path that reaches done, and stays through a return of
% prefixed_command itself.  Every kind of assignment; \setbox, where the box's own code takes it.  The return of a prefix before
% a non-prefixed command (WEB 1390): the token stays for the next assignment.  The returns of do_register_command, "You can't
% use x after \advance" (WEB 1415) and "Arithmetic overflow" (WEB 1414), are not that: control comes back and reaches done, so the
% token is inserted at once.  The goto done in WEB 1405 and in WEB 1430.  The token is saved once, the later one wins, and an
% undefined or a macro token behaves.
\input prelude
\errorcontextlines=0
\def\z{\message{[z]}}
\def\after#1{\afterassignment\z #1\message{[.]}}
\after{\def\a{}}
\after{\gdef\a{}}
\after{\edef\a{}}
\after{\xdef\a{}}
\after{\let\a=b}
\after{\let\a b}
\after{\futurelet\a\relax b}\par
\after{\chardef\a=1 }
\after{\mathchardef\a=1 }
\after{\countdef\a=1 }
\after{\dimendef\a=1 }
\after{\skipdef\a=1 }
\after{\muskipdef\a=1 }
\after{\toksdef\a=1 }
\after{\countdef\a=300 }
\after{\toks1={}}
\after{\toks1=\toks2 }
\after{\toks1=\toks300 }
\after{\everypar={}}
\after{\output={}}
\after{\toks300={}}
\after{\toks300=\toks1 }
\after{\count1=1 }
\after{\dimen1=1pt }
\after{\skip1=1pt }
\after{\muskip1=1mu }
\after{\tolerance=1 }
\after{\hsize=1pt }
\after{\catcode`a=11 }
\after{\mathcode`a=1 }
\after{\delcode`a=1 }
\after{\textfont1=\nullfont}
\after{\advance\count1 by 1 }
\after{\multiply\count1 by 1 }
\after{\divide\count1 by 1 }
\after{\advance\count300 by 1 }
\after{\prevdepth=1pt }
\after{\prevgraf=1 }
\after{\pagegoal=1pt }
\after{\deadcycles=1 }
\after{\interactionmode=1 }
\after{\wd0=1pt }
\after{\parshape=0 }
\after{\interlinepenalties=0 }
\after{\hyphenation{}}
\after{\patterns{}}
\after{\fontdimen1\tenrm=1pt }
\after{\hyphenchar\tenrm=1 }
\after{\lpcode\tenrm65=1 }
\after{\pdfnoligatures\nullfont}
\after{\font\a=cmr10 }
\after{\tenrm}
\after{\scrollmode}
\after{\nonstopmode}
\after{\openin3=prelude \read3 to\a}
\after{\global\def\a{}}
\after{\long\outer\def\a{}}
% the token comes after the assignment's own trailing input
\afterassignment\z\count1=1\relax
\afterassignment\z\count1=1 \relax
\afterassignment\relax\count1=1\show\z
\afterassignment x\count1=1\par
\afterassignment\undefinedone\count1=1
\afterassignment{\count1=1}
\afterassignment}\count1=1
\afterassignment\z\afterassignment\relax\count1=1 \message{[.]}
\afterassignment\relax\afterassignment\z\count1=1 \message{[.]}
\afterassignment\z\relax\relax\message{[not yet]}\count1=1 \message{[.]}
% the returns
\afterassignment\z\global\hrule \message{[not yet]}\count1=1 \message{[now]}
\afterassignment\z\advance x\par\message{[not yet]}\count1=1 \message{[now]}
\afterassignment\z\multiply\relax \message{[not yet]}\count1=1 \message{[now]}
\afterassignment\z\divide\count1 by 0 \message{[not yet]}\count1=1 \message{[now]}
\afterassignment\z\multiply\skip1 by 0 \message{[reaches done]}
\count1=2147483647 \afterassignment\z\multiply\count1 by 2 \message{[not yet]}\count1=1 \message{[now]}
\skip1=1pt \afterassignment\z\divide\skip1 by 0 \message{[not yet]}\count1=1 \message{[now]}
% \setbox: the token goes inside the box
\afterassignment\z\setbox0=\hbox{\message{[in box]}}\message{[.]}
\afterassignment\z\setbox0=\box2 \message{[.]}
\afterassignment\z\setbox0=x\par \message{[.]}
\afterassignment\z\setbox0=\vtop{\message{[in box]}}\message{[.]}
% in \accent's and the display's assignments
\tenrm\afterassignment\z\accent65 \count1=1 a\par
\accent65 \afterassignment\z\count1=1 \message{[.]}a\par
$$\halign{#\cr a\cr}\afterassignment\z\count1=1 $$
\end
